\documentclass{article}
\usepackage[utf8]{inputenc}
\usepackage{a4wide}

\begin{document}

\section*{หาวันว่าง}

มีเพื่อนกลุ่มหนึ่งจำนวน M คน ข้อมูลกำหนดช่วงวันที่แต่ละคนต้องทำงาน

(เป็นวันเริ่มและวันสิ้นสุดของแต่ละช่วง) โดยที่ช่วงต่างๆ นี้เรียงลำดับตามวันและไม่ซ้อนทับกัน

ให้หาช่วงวันทุกช่วงที่เพื่อนกลุ่มนี้ทุกคนว่างหมดและเรียงตามลำดับวันด้วย



\textbf{หมายเหตุ:}



\begin{enumerate}

\tightlist

\item

  จะไม่นับช่วงวันว่างที่ความยาวเป็น 0 เช่น 5 5 หรือ 13 13

\item

  หากไม่พบช่วงวันที่ทุกคนว่างได้เลย ให้ตอบว่า -1 -1

\end{enumerate}



\textbf{รูปแบบข้อมูลเข้า}



บรรทัดแรก จำนวนเต็มหนึ่งค่า M กำหนดจำนวนสมาชิกในกลุ่มเพื่อน M Î {[}2...50{]}



บรรทัดที่ 2 ถึง 2 + M จำนวนเต็มเป็นค่า N{[}i{]} โดยที่ i Î {[} 1... M {]}.

N{[}i{]} คือจำนวนช่วงทำงานของเพื่อนคนที่ i แล้วตามด้วยคู่ลำดับ (วันเริ่มและวันสิ้นสุด)

ของแต่ละช่วงวันทำงาน เว้นด้วยช่องว่าง 1 ช่อง



\textbf{รูปแบบข้อมูลออก}



คู่ลำดับของช่วงวันที่ทุกคนว่าง เรียงตามลำดับ เว้นด้วยช่วงว่าง



{\def\LTcaptype{none} % do not increment counter

\begin{longtable}[]{@{}

  >{\raggedright\arraybackslash}p{(\linewidth - 4\tabcolsep) * \real{0.3333}}

  >{\raggedright\arraybackslash}p{(\linewidth - 4\tabcolsep) * \real{0.3333}}

  >{\raggedright\arraybackslash}p{(\linewidth - 4\tabcolsep) * \real{0.3333}}@{}}

\toprule\noalign{}

\begin{minipage}[b]{\linewidth}\centering

Input

\end{minipage} & \begin{minipage}[b]{\linewidth}\centering

Output

\end{minipage} & \begin{minipage}[b]{\linewidth}\centering

คำอธิบาย

\end{minipage} \\

\midrule\noalign{}

\endhead

\bottomrule\noalign{}

\endlastfoot

\begin{minipage}[t]{\linewidth}\raggedright

3\\

2 1 2 5 6\\

1 1 3\\

1 4 10\\

\strut

\end{minipage} & \begin{minipage}[t]{\linewidth}\raggedright

3 4\\

\strut

\end{minipage} & \begin{minipage}[t]{\linewidth}\raggedright

มี 3 คน คนแรกมีช่วงทำงาน 2 ช่วง คือ {[}1, 2{]} และ {[}5, 6{]}\\

คนที่ 2 มีช่วงทำงาน 1 ช่วง คือ {[}1, 3{]}\\

คนที่ 3 มีช่วงทำงาน 1 ช่วง คือ {[}4, 10{]}\\

พบว่าทั้งสามคนว่างช่วงเวลา {[}3, 4{]}\\

\strut

\end{minipage} \\

\begin{minipage}[t]{\linewidth}\raggedright

3\\

2 1 3 6 7\\

1 2 4\\

2 2 5 9 12\\

\strut

\end{minipage} & \begin{minipage}[t]{\linewidth}\raggedright

5 6 7 9\\

\strut

\end{minipage} & \begin{minipage}[t]{\linewidth}\raggedright

มี 3 คน คนแรกมีช่วงทำงาน 2 ช่วง คือ {[}1, 3{]} และ {[}6, 7{]}\\

คนที่ 2 มีช่วงทำงาน 1 ช่วง คือ {[}2, 4{]}\\

คนที่ 3 มีช่วงทำงาน 2 ช่วง คือ {[}2, 5{]} และ {[}9, 12{]}\\

พบว่าทั้งสามคนว่างช่วงเวลา {[}5, 6{]} และ {[}7, 9{]}\\

\strut

\end{minipage} \\

\end{longtable}

}



\textbf{ข้อจำกัด}



จำนวนเพื่อนในกลุ่ม M ไม่เกิน 50



\subsection*{Input}
บรรทัดแรก จำนวนเต็มหนึ่งค่า M กำหนดจำนวนสมาชิกในกลุ่มเพื่อน M Î {[}2...50{]}



บรรทัดที่ 2 ถึง 2 + M จำนวนเต็มเป็นค่า N{[}i{]} โดยที่ i Î {[} 1... M {]}.

N{[}i{]} คือจำนวนช่วงทำงานของเพื่อนคนที่ i แล้วตามด้วยคู่ลำดับ (วันเริ่มและวันสิ้นสุด)

ของแต่ละช่วงวันทำงาน เว้นด้วยช่องว่าง 1 ช่อง



\subsection*{Output}
{คู่ลำดับของช่วงวันที่ทุกคนว่าง เรียงตามลำดับ เว้นด้วยช่วงว่าง}\\



\end{document}
